% Source-aligned English for printed p.160.
\SourcePage{160}

\LayerTarget{tohoku:remark:fine-sheaves-are-soft:p2}
\noindent of fine sheaves in terms of soft sheaves.

\LayerTarget{tohoku:construction:godement-resolution}
\textit{Resolution of the identity.} For every sheaf $F$, let $C^0(F)$ be the
product sheaf defined by the family of sets $F(x)$. There is a natural
injective homomorphism $F\longrightarrow C^0(F)$. Restricting to abelian
sheaves on $X$, the method of 2.5, Example a), constructs a resolution
$C(F)$ of the identity. It agrees with $C^0(F)$ in degree 0 and, in degree
$n\geq2$, is defined by
\[
C^n(F)=C^0\bigl(C^{n-1}(F)/\operatorname{Im}C^{n-2}(F)\bigr).
\]
The $C^n(F)$ are flabby and hence $\Gamma_\Phi$-acyclic for every antifilter
$\Phi$ of closed subsets. Thus $\Gamma_\Phi C(F)$ is a resolving functor for
$\Gamma_\Phi$, and
\[
H^n_\Phi(X,F)=H^n\bigl(\Gamma_\Phi C(F)\bigr).
\]
$C(F)$ is called the \textit{canonical resolution of $F$}; it was introduced
and used systematically by Godement.

\LayerTarget{tohoku:construction:cartan-resolution}
If $\Phi$ is paracompactifying, another resolving functor for $\Gamma_\Phi$
is obtained by taking a fixed resolution $C$ of the constant sheaf
$\mathbf Z$ by fine torsion-free sheaves and assigning to every $F$ the
complex of sheaves $F\otimes C$. This is a resolution of $F$, since $C$ is
torsion-free, and it is exact in $F$ for the same reason. Moreover, the
$F\otimes C^n$ are fine and hence $\Gamma_\Phi$-acyclic. Therefore
$\Gamma_\Phi(F\otimes C)$ is a resolving functor for $\Gamma_\Phi$, and
\[
H^n_\Phi(X,F)=H^n\bigl(\Gamma_\Phi(F\otimes C)\bigr).
\]
$F\otimes C$ will be called the Cartan resolution of $F$. Recall that it can
be used only when $\Phi$ is paracompactifying.

\LayerTarget{tohoku:example:irreducible-constant-sheaves}
\textit{An amusing example.} A space $X$ is called \textit{irreducible} if it
is not the union of two closed proper subsets, equivalently if every two
nonempty open subsets have nonempty intersection. This is also equivalent to
every open subset of $X$ being connected. Every \textit{constant} sheaf $F$
on $X$ is then plainly flabby. The converse also holds when $X$ is connected,
as one sees by taking a constant sheaf whose fiber has more than one point.
\textit{In particular, if $F$ is a constant sheaf of abelian groups on an
irreducible space $X$, then $H^p(X,F)=0$ for $p>0$.}

\LayerTarget{tohoku:section:3.4}
\textbf{3.4. Applications to lifting problems for structure
groups\textsuperscript{8)}.} Let $X$ be an irreducible algebraic variety over
an algebraically closed field $k$; we follow the terminology of [15,
Chapter II]. Let $\mathbf O$ be its sheaf of local rings, that is, the sheaf
of germs of regular functions on $X$, and let $\mathbf K$ be the sheaf of
germs of rational functions on $X$. The sheaf $\mathbf K$ is constant; see
the cited work, Proposition 9. Let $\mathbf O^*$ and $\mathbf K^*$ be the
subsheaves of $\mathbf O$ and $\mathbf K$, respectively, formed by germs of
invertible functions. Evidently, $\mathbf K^*$ is again a constant sheaf of
abelian groups, and $\mathbf O^*$ is a subsheaf of it.

\LayerTarget{tohoku:construction:divisor-sequence:p1}
The quotient sheaf $\mathbf K^*/\mathbf O^*=\mathbf D$ is the sheaf of
\textit{germs of locally principal divisors on $X$}. It agrees with the sheaf
of \textit{germs of divisors on $X$} if the local rings $\mathbf O(x)$ of
$X$ are factorial, for example if $X$ is nonsingular; we assume this from
now on. On the other hand, it is immediate that the sheaf $\mathbf D$ of
germs of divisors on

\LayerTarget{tohoku:footnote:8}
\begin{quote}
\footnotesize
8) This section is not needed for understanding what follows.
\end{quote}

\TohokuContinuation{161}

